% !TeX root = ../main.tex

\section{Métodos de análisis de resultados}
Para evaluar objetivamente los resultados de los algoritmos de deconvolución, se emplean dos parámetros métricos que permiten cuantificar la calidad de la reconstrucción utilizando datos sintéticos, dado que en estos casos se dispone de la señal original para su comparación directa. 

\subsection{Raíz del error cuadrático medio (RMSE)}
Este parámetro mide la magnitud promedio del error entre la señal original (ideal o de referencia) y otra señal distinta, ya sea la propia señal emborronada por el instrumento o la señal reconstruida tras aplicar un algoritmo de deconvolución. La ecuación que define esta métrica es la siguiente:

\begin{equation}
    \label{eq:RMSE}
    RMSE = \sqrt{\frac{1}{N}\sum_{i=1}^{N}(x_i - y_i)^2}
\end{equation}

donde $x_i$ es el valor de la señal original en el punto (o píxel) $i$, $y_i$ es el valor de la señal a comparar en ese mismo punto, y $N$ es el número total de elementos discretos de la señal. 

Dado que penaliza las diferencias elevándolas al cuadrado, un valor de RMSE más bajo indica que la señal recuperada es cuantitativamente más fiel a la señal original. Por el contrario, un RMSE más alto evidencia una mayor discrepancia entre la señal procesada y la referencia.

\subsection{Índice de Similitud Estructural (SSIM)}
El SSIM es una métrica de calidad de imagen que, a diferencia del RMSE, evalúa la degradación de la estructura espacial. Considera tres factores fundamentales: luminancia, contraste y estructura. El resultado de esta métrica es un valor comprendido entre $-1$ y $1$, donde los valores cercanos a $1$ indican una alta similitud con la imagen original. 

Aunque matemáticamente pueden presentarse valores negativos correspondientes a una correlación estructural inversa, la degradación instrumental estudiada no produce tales efectos. Por tanto, el rango de valores de interés físico para este análisis se restringe a $[0,1]$. 

Dada la complejidad analítica en el cálculo explícito de esta métrica, se recurre a la implementación computacional provista por la función \texttt{ssim} de la biblioteca \texttt{scikit-image} en Python, facilitando la evaluación eficiente del índice entre pares de imágenes. 

Cabe destacar que, si bien el RMSE y el SSIM proporcionan evaluaciones cuantitativas rigurosas, la inspección visual cualitativa de las imágenes sigue siendo un criterio complementario fundamental para dictaminar la viabilidad final de la reconstrucción.
